\chapter{一次分数型と連立漸化式}\label{章：一次分数型と連立漸化式}
\begin{bookcolumns}
線形の連立漸化式から成分の比だけを取り出すと,一次分数型が現れる.
本章ではこの対応を出発点にして,不動点からの差の比を使う理由を説明する.
行列・不動点・複素平面での回転を,最後に一つの具体例で結び付けよう.
\section{成分の比を追う}\label{節：成分の比を追う}

一次分数型
\[
 a_{n+1}=\frac{pa_n+q}{ra_n+s}
\]
を考える. 本章では係数は実数,$r\ne0$,$ps-qr\ne0$とする.
まず二つの数の比として$a_n$を記録してみよう.
$(x_1,y_1)=(a_1,1)$から
\[
 \begin{pmatrix}x_{n+1}\\y_{n+1}\end{pmatrix}
 =M\begin{pmatrix}x_n\\y_n\end{pmatrix},
 \qquad
 M=\begin{pmatrix}p&q\\r&s\end{pmatrix}
\]
と更新すれば,
\[
 \frac{x_{n+1}}{y_{n+1}}
 =\cfrac{p\cfrac{x_n}{y_n}+q}{r\cfrac{x_n}{y_n}+s}.
\]
従って,\,分母が非零の範囲では$a_n=\frac{x_n}{y_n}$となる.
一次の連立更新を比として記録すると,一次分数型が生まれるのである.

二つの成分を同じ非零数で倍しても比は変わらない.
長さと矢印の正負を区別せず,\,原点を通る直線としての方向を記録するこの考え方を,\,\textbf{射影的な見方}という.
分母$y_n=0$は,\,成分比を有限の実数で記録できなくなる方向に当たる.
元の実数列で定義できる範囲と,\,行列で状態を追える範囲は区別する.

不動点$\alpha$については
\[
 M\begin{pmatrix}\alpha\\1\end{pmatrix}
 =(r\alpha+s)\begin{pmatrix}\alpha\\1\end{pmatrix}.
\]
従って,\,\textbf{不動点は固有値ではなく,\,固有ベクトルの成分比}である.
対応する固有値は$\lambda_\alpha=r\alpha+s$であり,\,役割が異なる.
\NewFigProjective


ただし固有ベクトルの成分比が分かっても,一般の初期値からの比の変化はまだ分からない.
次節では,二つの不動点を基準にすれば等比数列を作れることを確かめる.

\section{不動点からの差の比}\label{節：不動点からの差の比}
\subsection{商は長さの比と角度の差}

二点$\alpha,\beta$を基準にして,\,点$z$の位置を
\[
 w=\frac{z-\alpha}{z-\beta}\qquad(z\ne\beta)
\]
で記録してみよう.
絶対値は
\[
 |w|=\frac{|z-\alpha|}{|z-\beta|}
\]
だから,\,二つの基準点からの距離の比を表す.
$z\ne\alpha$なら偏角は,ベクトル$z-\beta$から$z-\alpha$への有向角である.
二つの矢印を同じ始点に平行移動して角度を比較してもよい.
$z=\alpha$では$w=0$となり,偏角は定めない.
\NewFigRatio

この比が一定倍率で変わるなら,\,距離の比と角度の差の変化を同時に追える.
この比を使って,一次分数型の更新を簡単にしよう.

\begin{bookpoint}
複素数の積は\textbf{倍率と回転角}をまとめて記録する.
複素数の商は\textbf{長さの比と角度の差}をまとめて記録する.
\end{bookpoint}

\subsection{二つの不動点を使う}

次の項が
\[
 a_{n+1}=f(a_n),\qquad
 f(z)=\frac{pz+q}{rz+s}
\]
によって決まる場合を考える.
この形の$f$を\textbf{一次分数変換}という.
本節では係数は実数,$r\ne0$,$ps-qr\ne0$とし,\,元の更新で分母が零にならない範囲を扱う.
$r=0$なら$a_{n+1}=Aa_n+B$の形へ戻り,$ps-qr=0$なら定義できる入力で右辺が定数になるため,\,別に調べればよい.

不動点$\alpha$は
\[
 \alpha=\frac{p\alpha+q}{r\alpha+s}
\]
を満たすので,
\[
 r\alpha^2+(s-p)\alpha-q=0
\]
から求まる.
この方程式の根では$r\alpha+s\ne0$である.
もし零なら,\,元の方程式から$p\alpha+q=0$にもなり,$ps-qr\ne0$に反する.

不動点の式を使って引き算すると
\[
 f(z)-\alpha
 =\frac{(p-r\alpha)(z-\alpha)}{rz+s}.
\]
定数項は消えるが,\,まだ分母に$z$が残る.
ここで,\,異なるもう一つの不動点$\beta$があれば
\[
 f(z)-\beta
 =\frac{(p-r\beta)(z-\beta)}{rz+s}.
\]
二式の比を取れば,\,共通の分母が消える.
\[
 \frac{f(z)-\alpha}{f(z)-\beta}
 =\frac{p-r\alpha}{p-r\beta}
  \frac{z-\alpha}{z-\beta}.
\]
つまり
\[
 b_n=\frac{a_n-\alpha}{a_n-\beta},\qquad
 \kappa=\frac{p-r\alpha}{p-r\beta}
\]
と置くと,$b_{n+1}=\kappa b_n$である.
ここで$p-r\alpha,\,p-r\beta$も非零である.
実際,\,不動点の式から
$(p-r\alpha)(r\alpha+s)=ps-qr$となる.

初項が不動点なら定数列として先に扱う.
$a_1\ne\beta$なら,\,上の差の式から$a_n\ne\beta$が続くため,\,比を取れる.
解いた後は
\[
 b_n=b_1\kappa^{n-1},\qquad
 a_n=\frac{\beta b_n-\alpha}{b_n-1}
\]
で元へ戻す.
$b_n=1$になる場合は有限の$a_n$へ戻せず,\,元の更新もその項まで実数列としては定義できない.

\begin{bookpoint}
一つの不動点を引くと定数項が消える.
二つの不動点からの差の比を取ると,\,残った共通の分母も消える.
置き換えは,\,消したい部分に合わせて組み立てられる.
\end{bookpoint}

\noindent\textbf{【例題】}\par\nobreak
$a_1=3$,$a_{n+1}=\frac{2a_n+1}{a_n+2}$を解け.

\noindent\textbf{【方針】}\par\nobreak
不動点1と$-1$からの差の比を取ります.

\noindent\textbf{【解答】}\par\nobreak
\[
 a_n=\frac{2\cdot3^{n-1}+1}{2\cdot3^{n-1}-1}.
\]
\noindent\textbf{【解法】}\par\nobreak
$b_n=\frac{a_n-1}{a_n+1}$とすると
$b_{n+1}=\frac13b_n$,$b_1=\frac12$である.
従って$b_n=\frac1{2\cdot3^{n-1}}$となり,
$a_n=\frac{1+b_n}{1-b_n}$から答えを得る.
分母は正で,\,各項も正なので元の分母$a_n+2$は零にならない.
また$b_n\to0$だから$a_n\to1$となる.
不動点からの距離そのものではなく,\,二つの差の比が縮小していた.

\subsection{不動点が重なる場合}

不動点の方程式が重解$\alpha$を持つと,\,異なる二つの差を作れない.
この場合は
\[
 2r\alpha=p-s,\qquad
 p-r\alpha=r\alpha+s=\ell
\]
となり,$\ell\ne0$である.
$u_n=a_n-\alpha$と置くと
\[
 u_{n+1}=\frac{\ell u_n}{ru_n+\ell}.
\]
初項が不動点の場合を除き,\,逆数を取って
\[
 \frac1{u_{n+1}}=\frac1{u_n}+\frac r\ell
\]
という等差型へ直せる.
二つの固有方向が分かれない場合に足し算が残った\sectionref{節：重解で現れる足し算}と同じように,等比的な変化だけでは捉えられない足し算が残った.

例えば
\[
 a_1=5,\qquad a_{n+1}=\frac{4a_n-9}{a_n-2}
\]
では$\alpha=3$であり,
\[
 \frac1{a_{n+1}-3}=\frac1{a_n-3}+1.
\]
従って$a_n=3+\frac2{2n-1}$となる.
この例では$a_n>3$が保たれ,\,元の分母も非零である.

\subsection{差の比は固有成分の比}

異なる二つの不動点$\alpha,\beta$に対し,\,状態を
\[
 \begin{pmatrix}x_n\\y_n\end{pmatrix}
 =c_n\begin{pmatrix}\alpha\\1\end{pmatrix}
 +d_n\begin{pmatrix}\beta\\1\end{pmatrix}
\]
と分ける. 固有値を$\lambda_\alpha,\lambda_\beta$とすれば
\[
 c_{n+1}=\lambda_\alpha c_n,\qquad
 d_{n+1}=\lambda_\beta d_n.
\]
成分を比較すると
\[
 x_n-\alpha y_n=(\beta-\alpha)d_n,
\]
\[
 x_n-\beta y_n=(\alpha-\beta)c_n.
\]
従って,\,比が定義できる場合には
\[
 \frac{a_n-\alpha}{a_n-\beta}=-\frac{d_n}{c_n}
\]
となり,\,その公比は$\frac{\lambda_\beta}{\lambda_\alpha}$である.
不動点の和$r(\alpha+\beta)=p-s$から
\[
 \lambda_\beta=p-r\alpha,\qquad
 \lambda_\alpha=p-r\beta
\]
なので,\,前の一次分数型の公比と一致する.

本節でも$r\ne0,\,ps-qr\ne0$を仮定する.
このとき$M$は逆行列を持ち,\,固有値も非零である.
不動点を初項とする定数列を除けば,\,必要な成分比を取り続けられる.
元の更新の分母が零にならない条件は別に確かめる.

実係数で不動点が共役な非実数なら,\,固有値も共役である.
両者は絶対値が等しいため,\,\textbf{比を取ると共通の拡大率が消えて回転だけが残る}.
\sectionref{節：隣接3項間漸化式と回転}では回転に拡大・縮小が伴った. 次節で扱う一次分数型の円周上の座標では,成分比を取ることで共通の倍率が消える.

不動点が重なる場合は,\,行列も二つの固有方向に分けられない.
$r\ne0$だから$M$は単位行列の定数倍ではなく,\,固有方向は一本になる.
その寄与の蓄積を比で読み直すと,\,逆数を取った等差型が現れる.
隣接3項間漸化式の重解と,一次分数型の重なった不動点は,同じ行列の性質から生じている.

\section{実軸と単位円}\label{節：実軸と単位円}
\subsection{実数の数列を単位円へ}

不動点の二解が非実数なら,\,実係数なので$\alpha,\overline\alpha$という共役の組になる.
実数$a_n$について
\[
 b_n=\frac{a_n-\alpha}{a_n-\overline\alpha}
\]
と置く.
分子と分母は共役だから,$|b_n|=1$である.
また
\[
 \kappa=\frac{p-r\alpha}{p-r\overline\alpha}
\]
も分子と分母が共役なので$|\kappa|=1$となる.
従って,$b_{n+1}=\kappa b_n$は\textbf{単位円周上の一定角の回転}である.
実軸上で複雑に見えた動きが,\,円周上では単純になる.

\noindent\textbf{【例題】}\par\nobreak
\[
 a_1=\frac12,\qquad a_{n+1}=\frac{a_n+1}{1-a_n}
\]
を複素平面で解釈し,\,周期を求めよ.

\noindent\textbf{【方針】}\par\nobreak
不動点は$i,-i$です. 差の比を取り,\,一回の回転角を調べます.

\noindent\textbf{【解答】}\par\nobreak
最小周期は4であり,
\[
 \frac12,\ 3,\ -2,\ -\frac13,\ \frac12,\ldots
\]
と繰り返す.

\noindent\textbf{【解法】}\par\nobreak
$b_n=\frac{a_n-i}{a_n+i}$とすると
\[
 b_{n+1}=\frac{1+i}{1-i}b_n=ib_n.
\]
毎回反時計回りに$\frac{\pi}{2}$回転する.
$b_n$は零でなく,4回より少ない正の回数では元へ戻らない.
復元は$a_n=i\frac{b_n+1}{1-b_n}$である.
実際に元の式で四項を計算すると上の並びを得る.
どの項も1ではないため,\,すべての更新が定義できる.

ここで,\,回転が定義できることと,\,元の実数列が全項で定義できることは区別する.
例えば同じ式で$a_1=0$とすると$a_2=1$となり,$a_3$は定義できない.
円周上では回転を続けられるが,\,復元式の分母が零になる点に到達している.

\subsection{実軸と円を結ぶ座標}\label{節：実軸と円を結ぶ座標}

回転を三角関数と結び付けるため,\,別の一定倍をした座標
\[
 w=\Phi(a)=\frac{1+ia}{1-ia}
\]
を使おう.
実数$a$なら$|w|=1$であり,\,この$\Phi(a)$は先ほどの$b$に対して$w=-b$である.
例えば$a=0$は$w=1$,$a=1$は$w=i$へ移る.

$a=\tan\theta$と書ける範囲では,
\[
 \begin{aligned}
 \Phi(\tan\theta)
 &=\frac{\cos\theta+i\sin\theta}{\cos\theta-i\sin\theta}\\
 &=\cos2\theta+i\sin2\theta.
 \end{aligned}
\]
実数$a$を傾きとして記録した角度が,\,円周上ではその2倍の角度になる.
\NewFigCayley

$w=u+iv$が単位円上で$w\ne-1$なら,\,逆に
\[
 a=\frac{v}{1+u}
\]
で戻れる.
点$(-1,0)$から$(u,v)$への直線の傾きが$a$である.
従って$\Phi$は実軸と,\,一点$-1$を除いた単位円を一対一に対応させる.
除いた一点は,\,傾きが無限に大きくなる方向に当たる.
必要に応じて実軸に無限遠点$\infty$を一つ加えると円全体と対応するが,\,実数の割り算では依然として分母の条件が必要である.

実数$t$を固定し,
\[
 f(a)=\frac{a+t}{1-ta}
\]
とすると,\,直接の計算から
\[
 \Phi(f(a))=\Phi(t)\Phi(a)
\]
となる. $t=\tan\phi$なら$\Phi(t)=e^{2i\phi}$であり,
元の式は$a=\tan\theta$の角度を$\theta+\phi$へ進める式である.
正接の加法公式
\[
 \tan(\theta+\phi)
 =\frac{\tan\theta+\tan\phi}{1-\tan\theta\tan\phi}
\]
が,\,円周上の掛け算として現れた.
角度は$\pi$の整数倍だけ違っても同じ傾きを表す.

\section{三つの見方を結ぶ}\label{節：三つの見方を結ぶ}
\subsection{一つの例を三通りに記録する}

特に
\[
 M=\begin{pmatrix}1&1\\-1&1\end{pmatrix}
\]
を使おう.
状態の更新は$x_{n+1}=x_n+y_n,\,y_{n+1}=-x_n+y_n$である.
複素数を$z_n=y_n+ix_n$と選ぶと,
\[
 z_{n+1}=(1+i)z_n.
\]
この座標では,\,毎回$\sqrt2$倍されながら反時計回りに$\frac{\pi}{4}$回転する.
各成分は特性方程式$X^2-2X+2=0$に対応し,
\[
 x_{n+2}=2x_{n+1}-2x_n
\]
などを満たす.

一方,\,成分比$a_n=\frac{x_n}{y_n}$を記録すると,
\[
 a_{n+1}=\frac{a_n+1}{1-a_n}.
\]
共通の拡大率が消え,\,傾きの変化だけが残る.
さらに$w_n=\Phi(a_n)$で円周へ移すと,
\[
 w_{n+1}=iw_n.
\]
角度を2倍して記録する$\Phi$によって,\,回転角は$\frac{\pi}{2}$になる.

$(x_1,y_1)=(\frac12,1)$なら,\,成分比は先ほどの4周期の列となる.
行列では4回後の状態は$M^4\boldsymbol v_1=-4\boldsymbol v_1$であり,\,元の状態そのものには戻らない.
しかし両成分が同じ倍率を受けるため比は戻る.
\textbf{状態の周期と,\,比の周期は同じとは限らない}.


\begin{bookpoint}
同じ行列から,連立漸化式,隣接3項間漸化式,一次分数型を作れる.
状態全体を追う場合と,成分の一つや比だけを追う場合とでは,読み取れる性質が異なる.
\end{bookpoint}

\end{bookcolumns}
